% =====================================================================
% LilyTorch — Supplementary Information
%
% Detailed numerical formulations that are summarised only briefly in
% the main manuscript (main.tex):
%   S1  Spatial discretisation and ghost cells
%   S2  Advection schemes (the six convective limiters)
%   S3  Boundary conditions (velocity ghost layer + Poisson)
%   S4  Pressure-projection / Poisson back-ends
%   S5  Strong (implicit) fluid--structure coupling
%
% Standalone document: compiles on its own and shares references.bib
% with the main manuscript.
%
%   pdflatex supplementary
%   bibtex   supplementary
%   pdflatex supplementary
%   pdflatex supplementary
% =====================================================================

\documentclass[10pt,onecolumn]{article}

\usepackage[margin=1in]{geometry}
\usepackage{amsmath,amssymb,amsfonts}
\usepackage{graphicx}
\usepackage{cite}
\usepackage{booktabs}
\usepackage{array}
\usepackage{url}
\usepackage[hidelinks]{hyperref}
\usepackage{bm}
\usepackage{xspace}

% ---- Convenience macros (mirror main.tex) ---------------------------
\newcommand{\lilytorch}{\textsc{LilyTorch}\xspace}
\newcommand{\farms}{\textsc{FARMS}\xspace}
\newcommand{\mujoco}{\textsc{MuJoCo}\xspace}
\newcommand{\pytorch}{\textsc{PyTorch}\xspace}
\newcommand{\waterlily}{\textsc{WaterLily.jl}\xspace}
\newcommand{\bdim}{BDIM\xspace}

% Supplementary numbering: S1, S2, ... and equations (S1), (S2), ...
\renewcommand{\thesection}{S\arabic{section}}
\renewcommand{\theequation}{S\arabic{equation}}
\renewcommand{\thetable}{S\arabic{table}}
\renewcommand{\thefigure}{S\arabic{figure}}

\title{\lilytorch: Supplementary Information\\[2pt]
\large Detailed Numerical Formulations for the Advection Schemes,
Boundary Conditions, Pressure Solvers, and Implicit Coupling}
\author{Andrea Ferrario}
\date{}

\begin{document}
\maketitle

\noindent
This supplement expands the numerical methods that are stated only
briefly in the main manuscript.  Section and equation numbers carry an
``S'' prefix; cross-references of the form ``\S2.x'' or ``Eq.~(x)''
without the prefix point to the main text.  Symbols follow the main
text: $h$ is the (uniform) cell size, $\Delta t$ the time step,
$\mathbf u$ the velocity, $p$ the pressure, $\rho$ the density, $\nu$
the kinematic viscosity, and $d(\mathbf x)$ the signed-distance
function (SDF) of an immersed body (positive in the fluid).  All fields
live on a Marker-and-Cell (MAC) staggered grid with one ghost cell per
side.


% =====================================================================
\section{Spatial discretisation and ghost cells}\label{sec:s-grid}
% =====================================================================

The domain $\Omega\subset\mathbb R^{d}$ ($d\in\{2,3\}$) is stored as
arrays of shape $(N_x{+}2)\times(N_y{+}2)\,[\times(N_z{+}2)]$, where the
single extra layer on each side is the \emph{ghost} (halo) layer used to
impose boundary conditions.  Pressure $p$ lives at cell centres; the
velocity components live on their own staggered face grids ($u$ on
$x$-faces, $v$ on $y$-faces, $w$ on $z$-faces).  Throughout this
supplement we use the index convention that, along a given dimension,
slice \texttt{[0]} is the low ghost layer, \texttt{[-1]} the high ghost
layer, and \texttt{[1:-1]} the interior.

The convective operator acts on a one-dimensional family of stencils:
for a transported quantity $\phi$ and a face velocity
$f$, the flux through an interior face uses the three aligned values
$(\phi_{\text{upstream}},\phi_{\text{centre}},\phi_{\text{downstream}})$,
selected by the sign of $f$.  The transverse face velocity needed to
advect a staggered momentum component is obtained by averaging the
appropriate neighbouring velocities to the flux face; for a
cell-centred scalar transported along direction $d$ the $d$-face already
coincides with the MAC location of $u_d$, so no transverse averaging is
needed.  These details are what make the scheme functions below pure,
dimension-agnostic kernels $\lambda(u,c,d)$ of the three upstream/
centre/downstream samples.


% =====================================================================
\section{Advection schemes}\label{sec:s-advection}
% =====================================================================

The convective term is written in flux form,
$\nabla\!\cdot(\mathbf u\otimes\phi)$, and discretised by reconstructing
a face value $\phi_f$ from the three aligned cell values
$u=\phi_{\text{upstream}}$, $c=\phi_{\text{centre}}$,
$d=\phi_{\text{downstream}}$ (relative to the local flow direction).
Five of the six schemes return a single face value $\phi_f=\lambda(u,c,d)$
and differ only in the limiter $\lambda$; the sixth (semi-Lagrangian)
replaces the flux form entirely.  The advection--diffusion sub-step is
advanced with forward Euler and embedded in the outer Heun (RK2)
predictor--corrector of the main text (\S2.2).

Several limiters are conveniently expressed through the upwind ratio of
consecutive gradients
\begin{equation}
  r = \frac{c-u}{d-c},
  \label{eq:s-rf}
\end{equation}
evaluated with a small floor on the denominator; where $|d-c|$
underflows, the scheme falls back to the first-order upwind value $c$.
The TVD limiters then take the form
\begin{equation}
  \phi_f = c + \tfrac12\,(d-c)\,\psi(r),
  \label{eq:s-tvd}
\end{equation}
so the limiter function $\psi(r)$ fully characterises the scheme.

\paragraph{QUICK (third-order upwind).}
The Quadratic Upstream Interpolation for Convective Kinematics
scheme~\cite{leonard1979quick} fits a parabola through the three points
and is stabilised with a double-median (ULTIMATE-style) limiter:
\begin{align}
  \phi_f^{\text{QUICK}} &= \tfrac{1}{6}\,(5c + 2d - u),
  \label{eq:s-quick-raw}\\
  m_{\text{in}} &= \operatorname{median}(10c-9u,\; c,\; d),
  \nonumber\\
  \phi_f &= \operatorname{median}\!\bigl(\phi_f^{\text{QUICK}},\; c,\;
            m_{\text{in}}\bigr).
  \label{eq:s-quick}
\end{align}
The inner median bounds the curvature correction, and the outer median
clips the result into the monotone interval, which suppresses the
dispersive overshoots of the unlimited parabola while retaining
third-order accuracy in smooth regions.  This is the default scheme.

\paragraph{ADBQUICKEST (Courant-dependent TVD QUICKEST).}
A TVD form of Leonard's QUICKEST in which the limiter depends explicitly
on the local Courant number $C=|f|\,\Delta t/h$.  With $r$ from
\eqref{eq:s-rf} the limiter is
\begin{equation}
  \psi(r) = \max\!\Bigl[\,0,\;
    \min\!\bigl(\,2(1-C)\,r,\;
      \tfrac{(2+C^2-3C)+(1-C^2)\,r}{3(1-C)},\;
      2(1-C)\bigr)\Bigr],
  \label{eq:s-abdq}
\end{equation}
inserted into \eqref{eq:s-tvd}.  The explicit $C$-dependence tightens
the stability bound at large Courant numbers, so this scheme is the
preferred high-order option when $\Delta t$ approaches the advective CFL
limit.  The implementation requires the actual per-face Courant number
to be consistent with the $C$ used in the limiter (a global
$\max|f|\,\Delta t/h$ is used).

\paragraph{CUBISTA (Convergent and Universally Bounded Interpolation
Scheme).}
A second-order, universally bounded high-resolution scheme of Alves,
Oliveira \& Pinho~\cite{alves2003convergent}, designed for smooth
convergence of iterative solvers.  Its limiter is the piecewise-linear
envelope
\begin{equation}
  \psi(r) = \max\!\Bigl[\,0,\;
    \min\!\bigl(\,\tfrac{3}{2}\,r,\;
      \tfrac{3}{4}\,r + \tfrac14,\;
      \tfrac{3}{2}\bigr)\Bigr].
  \label{eq:s-cubista}
\end{equation}

\paragraph{van Leer.}
The classical smooth (differentiable) symmetric limiter~\cite{vanleer1979towards}
\begin{equation}
  \psi(r) = \frac{r + |r|}{1 + |r|},
  \label{eq:s-vanleer}
\end{equation}
again used in \eqref{eq:s-tvd}.  It is second-order and TVD, with a
gentler shape than the QUICK family.

\paragraph{Central differences (CDS).}
The plain non-upwind, non-TVD second-order reconstruction
\begin{equation}
  \phi_f = \tfrac12\,(c + d).
  \label{eq:s-cds}
\end{equation}
It is the most accurate in smooth, well-resolved flow but admits
dispersive wiggles when the cell Reynolds/P\'eclet number is large; it
is mainly used for verification and for the boundary-face fallback of
the limited schemes (the lowest/highest interior faces, where the full
three-point stencil is unavailable, default to CDS).

\paragraph{Semi-Lagrangian back-tracing.}
For very large time steps the unconditionally stable semi-Lagrangian
advection of Stam~\cite{stam1999stable} is available.  Instead of a flux
balance, the departure point $\mathbf x_d = \mathbf x - \Delta t\,
\mathbf u(\mathbf x)$ of each grid node is traced backwards and the field
is interpolated there,
\begin{equation}
  \phi^{*}(\mathbf x) = \phi^{n}\!\bigl(\mathbf x - \Delta t\,
  \mathbf u^{n}(\mathbf x)\bigr),
  \label{eq:s-semilag}
\end{equation}
by trilinear (3-D) / bilinear (2-D) interpolation.  This removes the
advective CFL restriction entirely at the cost of substantial numerical
diffusion, and is intended for fast, low-fidelity previews rather than
quantitative runs.

\medskip
Table~\ref{tab:s-advection} summarises the schemes.  All five flux-form
limiters share the registry used both by the momentum advection and by
the volume-of-fluid colour transport of the two-phase model
(main text \S2.7); they are additionally available as a single fused
CUDA flux kernel that accumulates the right-hand side in place
(main text \S2.10).

\begin{table}[t]
\centering
\caption{Convective schemes exposed through the
\texttt{convection\_method} configuration key.}
\label{tab:s-advection}
\begin{tabular}{@{}l l c l@{}}
\toprule
Key & Scheme & Order & Notes\\
\midrule
\texttt{quick}        & QUICK + double median   & 3 & default; ULTIMATE-limited \\
\texttt{abdquickest}  & ADBQUICKEST             & 3 & TVD, Courant-dependent \\
\texttt{cubista}      & CUBISTA                 & 2 & TVD, universally bounded \\
\texttt{van\_leer}    & van Leer                & 2 & TVD, smooth limiter \\
\texttt{cds}          & central differences     & 2 & not TVD; verification / fallback \\
\texttt{semi-lagrangian} & Stam back-tracing    & 1--2 & unconditionally stable, diffusive \\
\bottomrule
\end{tabular}
\end{table}


% =====================================================================
\section{Boundary conditions}\label{sec:s-bc}
% =====================================================================

\lilytorch distinguishes the \emph{velocity} boundary conditions, which
act on the ghost layer of each MAC component, from the \emph{Poisson}
boundary condition, which is a property of the chosen pressure back-end.

\subsection{Velocity ghost-cell conditions}

Each velocity component carries an independent list of per-face
condition types and values.  The faces are ordered
\begin{equation*}
  (\text{dim0\_lo},\ \text{dim0\_hi},\ \text{dim1\_lo},\ \text{dim1\_hi},
   \ [\text{dim2\_lo},\ \text{dim2\_hi}]),
\end{equation*}
i.e.\ $(\text{west, east, south, north}[, \text{bottom, top}])$, giving
a list of $4$ entries in 2-D and $6$ in 3-D per component
(\texttt{bc\_type\_u/v/w} and \texttt{bc\_values\_u/v/w} in the
configuration; an unspecified tail defaults to Neumann).  Three
elementary ghost operations cover all supported cases.  Writing $g$ for
the ghost slice and $i$ for the adjacent interior slice:

\begin{itemize}
\item \textbf{Neumann (zero-gradient)} — copy the interior value into
  the ghost,
  \begin{equation}
    \phi_g = \phi_i,
    \label{eq:s-neumann}
  \end{equation}
  giving $\partial\phi/\partial n = 0$ at the face.  This is the default
  open/outflow-like condition.
\item \textbf{Dirichlet (direct)} — for a velocity component
  \emph{normal} to its own boundary face the wall value coincides with a
  grid location, so it is written directly,
  \begin{equation}
    \phi_g = \phi_{\text{wall}}.
    \label{eq:s-dirichlet}
  \end{equation}
  This sets, e.g., a prescribed inflow normal velocity or a zero
  through-wall velocity.
\item \textbf{Dirichlet (reflective / mirror)} — for a component
  \emph{tangential} to the face the wall value sits half a cell outside
  the last interior node, so the prescribed value $\phi_{\text{wall}}$ is
  imposed by linear reflection,
  \begin{equation}
    \phi_g = 2\,\phi_{\text{wall}} - \phi_i,
    \label{eq:s-reflect}
  \end{equation}
  which makes the face-interpolated value equal $\phi_{\text{wall}}$
  exactly.  With $\phi_{\text{wall}}=0$ this is the standard no-slip
  mirror.
\end{itemize}

The three operations are applied in a fixed, stable order
(Neumann, then direct Dirichlet, then reflective Dirichlet) so that a
wall-face direct write is already committed before the reflective form
reads the up-to-date interior value.  On the GPU, when all components
are contiguous CUDA tensors of the same floating dtype, the entire set
of per-face slice copies is dispatched as a single fused
\texttt{apply\_bcs\_2d}/\texttt{apply\_bcs\_3d} CUDA kernel (one launch
instead of many small slice assignments); an eager \pytorch loop is the
fallback otherwise.  A typical channel uses a direct Dirichlet inflow on
the west face, Neumann outflow on the east face, and reflective no-slip
on the lateral walls; a free-stream / open box uses Neumann on all
faces.

\subsection{Pressure (Poisson) boundary conditions}

The pressure-projection back-end carries its own boundary treatment,
selected by \texttt{poisson\_bc\_type}:

\begin{itemize}
\item \textbf{Neumann} ($\partial p/\partial n = 0$ on every face) — the
  homogeneous-Neumann Laplacian is exactly diagonalised by the type-II
  DCT (Section~\ref{sec:s-poisson}), giving the fastest solve.  Because
  the all-Neumann problem fixes $p$ only up to an additive constant, the
  pressure datum (null space) must be pinned; this is the natural choice
  for closed/walled domains.
\item \textbf{Free} (free-space / unbounded) — the right-hand side is
  zero-padded and convolved with a regularised free-space Green's
  function, emulating an infinite domain with no reflected pressure
  waves at the box edges.  This is the appropriate choice for an
  isolated body in open water, often combined with the sponge layer
  (main text \S2.6) on the velocity field.
\end{itemize}

The variable-coefficient multigrid/MGCG back-end (used for the BDIM
mask, variable density, and two-phase flows) inherits the homogeneous
Neumann condition at domain faces while embedding the immersed-body
condition through the face coefficients $c_{d\pm}$ rather than through
the ghost layer; see Section~\ref{sec:s-poisson}.  Note that the
velocity and pressure conditions are chosen consistently: a Dirichlet
(prescribed-velocity) inflow corresponds to a homogeneous-Neumann
pressure face, while a Neumann (zero-gradient) velocity outflow pairs
with a pressure datum constraint to keep the discrete projection
well-posed.


% =====================================================================
\section{Pressure-projection / Poisson solvers}\label{sec:s-poisson}
% =====================================================================

After the provisional velocity $\tilde{\mathbf u}$ is formed, the
incompressibility constraint is enforced by solving a Poisson problem
for $p$ and correcting
$\mathbf u^{n+1}=\tilde{\mathbf u}-(\,w\Delta t/\rho)\,\nabla p$.
\lilytorch ships three back-ends; the appropriate one depends on the
boundary condition and on whether the operator is constant- or
variable-coefficient.

\subsection{Spectral DCT solver (constant-coefficient Neumann)}

For homogeneous Neumann boundaries with a constant coefficient the
discrete Laplacian is diagonal in the type-II Discrete Cosine
Transform~\cite{makhoul1980dct}.  Each 1-D second difference has
eigenvalues
\begin{equation}
  \lambda_k = \frac{2}{h^2}\Bigl[\cos\!\Bigl(\tfrac{\pi k}{N}\Bigr)-1\Bigr],
  \qquad k=0,\dots,N-1,
  \label{eq:s-dct}
\end{equation}
and the $d$-dimensional eigenvalue is the sum over axes.  The solve is
therefore \emph{forward DCT $\rightarrow$ divide by
$\sum_{\text{axes}}\lambda$ $\rightarrow$ inverse DCT}, with the
zero-frequency (null-space) mode set to zero to fix the pressure datum.
This $\mathcal O(N\log N)$ path is the fastest and is used for the
majority of single-phase validation benchmarks.

\subsection{Free-space Green's function (unbounded)}

For an unbounded domain the right-hand side is zero-padded to $2N$ per
axis and convolved (via FFT) with a regularised free-space Green's
function, so the periodic wrap-around of the padded FFT does not
contaminate the interior.  In 2-D the eighth-order algebraic smoothing
of Hejlesen et al.~\cite{hejlesen2013high} is used; in 3-D the
Gaussian/erf regularisation
\begin{equation}
  G(r) = -\frac{1}{4\pi r}\,\operatorname{erf}\!\Bigl(\frac{r}{\sqrt 2\,
  \sigma}\Bigr),\qquad \sigma = h,
  \label{eq:s-green}
\end{equation}
removes the $1/r$ singularity at the origin.  The discrete Green's
kernel is precomputed once and cached on disk, so the per-step cost is
two padded FFTs plus a pointwise multiply.

\subsection{Geometric multigrid and MGCG (variable-coefficient)}

When the BDIM mask $\mu_0$ or a variable fluid/body density makes the
projection genuinely variable-coefficient, the discrete operator along
dimension $d$ is
\begin{equation}
  \frac{c_{d+}\,p_{i+1}-(c_{d+}+c_{d-})\,p_i+c_{d-}\,p_{i-1}}{h^2},
  \label{eq:s-mg-stencil}
\end{equation}
where the face coefficients $c_{d\pm}$ are the \emph{harmonic} average of
the cell values straddling each MAC face (of $\mu_0$ in the BDIM case,
or of $\Delta t/\rho$ in the variable-density / two-phase case).  The
harmonic average keeps the discrete Laplacian symmetric
positive-definite and well-balanced across the $10^3{:}1$ density jump
at an air--water interface.

The system is solved by a geometric V-cycle~\cite{trottenberg2001multigrid}:

\begin{itemize}
\item \textbf{Smoothers}: weighted Jacobi (default relaxation
  $\omega=0.7$) or red/black Gauss--Seidel.
\item \textbf{Restriction}: full-weighting of the residual to the
  coarse grid.
\item \textbf{Prolongation}: bilinear (2-D) / trilinear (3-D)
  interpolation of the coarse correction.
\item \textbf{Coarsest level}: a trivial $2\times2\,(\times2)$ direct
  solve.
\end{itemize}

For anisotropic production grids the V-cycle alone smooths weakly, so it
is optionally wrapped in a conjugate-gradient outer iteration (MGCG),
with the V-cycle acting as the preconditioner.  This roughly halves the
solve time on the strongly anisotropic two-phase grids (main text
\S2.10).  Two further accelerations are available and are bit-compatible
with the plain multigrid result: a \emph{recycled-Krylov} (deflated
MGCG) variant that reuses approximate near-null-space vectors across
time steps to help where the smoother stalls, and, for small grids
dominated by kernel-launch overhead, capture of the entire native solve
as a single CUDA graph.

\subsection{Variable-density and two-phase coupling}

When the body density $\rho_b$ differs from the fluid density $\rho_f$
(MuJoCo coupling) an effective density
\begin{equation}
  \rho(\mathbf x) = \rho_b + (\rho_f-\rho_b)\,\mu_0(\mathbf x)
  \label{eq:s-rho-eff}
\end{equation}
is assembled and the face coefficients $c=\Delta t/\rho_{\text{face}}$
are formed directly on the MAC faces.  The same machinery, with
$\rho(\alpha)=\alpha\rho_w+(1-\alpha)\rho_a$ from the volume-of-fluid
colour field $\alpha$, drives the two-phase free-surface model.  Because
the spectral back-ends assume a constant coefficient, all
variable-density and two-phase flows are solved exclusively by the
multigrid/MGCG path.


% =====================================================================
\section{Strong (implicit) fluid--structure coupling}\label{sec:s-implicit}
% =====================================================================

\emph{This section documents capability that is present in the solver
but only alluded to as future work in the main manuscript (limitations,
\S3.5).}  By default the fluid and \mujoco are advanced in a single
\emph{explicit}, loosely-coupled (staggered) pass per time step: the
fluid reads the body pose, integrates the hydrodynamic wrench, and
writes it back once.  For light or force-driven bodies whose
mean density approaches that of the fluid, this staggered scheme suffers
the well-known \emph{added-mass instability}~\cite{causin2005added}: the
explicit wrench lags the body acceleration it induces, and the coupling
diverges regardless of how small $\Delta t$ is.

\lilytorch optionally closes the coupling implicitly within each time
step, in the spirit of partitioned solvers such as
preCICE~\cite{bungartz2016precice}.  Let $x$ collect the interface
unknowns (the per-link wrench / interface state) passed from the body
solver to the fluid, and let $\mathcal F(x)$ be the value returned after
one fluid$\rightarrow$body sweep.  A converged time step is a fixed point
$x^{*}=\mathcal F(x^{*})$, found by sub-iterating
\begin{equation}
  \tilde x_k = \mathcal F(x_k),\qquad
  r_k = \tilde x_k - x_k,
  \label{eq:s-fixedpoint}
\end{equation}
until $\lVert r_k\rVert$ falls below a tolerance (or a maximum sweep
count is reached).  Three accelerators control how $x_{k+1}$ is built
from the residual history.

\paragraph{Constant under-relaxation.}
The classical staggered relaxation
\begin{equation}
  x_{k+1} = x_k + \omega\,r_k,
  \label{eq:s-constant}
\end{equation}
with a fixed factor $\omega\in(0,1]$.  $\omega=1$ reproduces the plain
explicit step; a small $\omega$ buys stability but smears the true
force and converges slowly.  This is the baseline.

\paragraph{Aitken (Irons--Tuck) adaptive relaxation.}
The scalar relaxation factor is recomputed every sub-iteration from the
last two residuals~\cite{irons1969version},
\begin{equation}
  \omega_k = -\,\omega_{k-1}\,
  \frac{\langle r_{k-1},\,r_k-r_{k-1}\rangle}
       {\lVert r_k-r_{k-1}\rVert^{2}},
  \label{eq:s-aitken}
\end{equation}
clamped in magnitude to $\omega_{\max}$ and applied as in
\eqref{eq:s-constant}.  This typically converges several times faster
than a hand-tuned constant $\omega$ while remaining a scalar one-liner;
the converged factor is carried into the next time step as the initial
guess.

\paragraph{Interface Quasi-Newton (IQN-ILS).}
The most robust option is the Interface Quasi-Newton with Inverse
Jacobian from a Least-Squares model of
Degroote et al.~\cite{degroote2009performance}.  Differences of
successive residuals and interface states are stacked into matrices
\begin{equation}
  V_k = [\,\Delta r_{k},\,\Delta r_{k-1},\,\dots\,],\qquad
  W_k = [\,\Delta \tilde x_{k},\,\Delta \tilde x_{k-1},\,\dots\,],
  \label{eq:s-iqn-cols}
\end{equation}
and the Newton update $\Delta x = -J^{-1}r_k$ is approximated by solving
the small least-squares problem $V_k\,\alpha \approx -r_k$ (via a QR
factorisation) and forming $\Delta x = W_k\,\alpha + (-r_k - V_k\alpha)$.
Columns may be \emph{reused} from a configurable number of previous time
steps (the \texttt{reuse} window) to enrich the Jacobian model and cut
the number of sweeps.

\medskip
The scheme and accelerator are selected through a \texttt{body.coupling}
configuration block (\texttt{scheme: explicit\,|\,implicit};
\texttt{accelerator: iqn-ils\,|\,aitken\,|\,constant}; \texttt{reuse},
\texttt{tol}, \texttt{max\_iter}).  In implicit mode the constant
force-relaxation low-pass used by the explicit path is bypassed.  A
practical caveat applies to column reuse: a single diverged
quasi-Newton sweep can poison the reuse window and silently degrade the
method back to an explicit step, so for delicate free-swimmer cases the
Aitken accelerator (or \texttt{reuse: 0}) is the safer default.  Implicit
coupling is intended for light or force-driven bodies where added mass
dominates; position-controlled swimmers, whose kinematics are
prescribed, are advanced with the cheaper explicit path.


% =====================================================================
\bibliographystyle{IEEEtran}
\bibliography{references}
% =====================================================================

\end{document}
